\section{课程地图与 Transformer 基础：架构只是第一层}

这堂 V5 开场课不是把 Transformer 重新讲成一张 attention 公式。四位讲者把现代系统拆成十个相互连接的层次：表示与 attention、数据配方、推理时计算、偏好优化、agent loop、视觉与脑成像接口，以及部署后的持续适应。读完整堂课时应不断追问：某个方法改变的是参数、数据、推理过程、外部工具，还是模型与环境之间的闭环？


\lecturefigure{slide-001.jpg}{CS25 Transformers United V5：课程与讲座入口}{CS25 V5 official deck, p.~1}

\subsection{课程承诺：从机制到边界}

课程 takeaways 把 attention、跨域应用、frontier research 与 limitations 放在同一页。这个顺序很重要：理解模型怎样工作，不等于理解系统为什么成功；看到应用扩张，也不等于忽略成本、可靠性与持续学习缺口。讲义因此把每个章节都写成“机制—证据—边界—工程决策”四步。


\lecturefigure{slide-012.jpg}{课程希望学生掌握的机制、应用、前沿方法与开放问题}{CS25 V5 official deck, p.~12}

\teachervoice{00:01:20--00:01:58，讲者把整堂 overview 明确分成 pretraining/data、post-training、applications 与 weaknesses。课堂提示：这些不是四个独立 buzzword，而是模型生命周期的四个干预位置。}

\begin{importantbox}{现代 Transformer 系统的四层接口}
\begin{enumerate}
  \item \textbf{Representation/architecture}：token 怎样变成向量并相互作用；
  \item \textbf{Training data/objective}：模型从哪些分布和损失中学习；
  \item \textbf{Inference/post-training}：部署前后如何追加搜索、偏好与工具；
  \item \textbf{Environment/update}：系统怎样行动、接收反馈并持续改变。
\end{enumerate}
后文所有方法都应先归到其中一层，再讨论跨层耦合。
\end{importantbox}

\subsection{Embedding：离散 token 怎样进入连续空间}

文本不能直接作为矩阵运算输入。设词表大小为 $|V|$、隐藏维度为 $d$，embedding table 为 $E\in\mathbb{R}^{|V|\times d}$，第 $i$ 个 token id 为 $x_i$，则
\[
e_i=E[x_i].
\]
其中 $e_i\in\mathbb{R}^{d}$ 是 token 的稠密向量。静态 embedding 给同一词同一向量；contextual embedding 则在后续层根据上下文更新表示，因此 ``bank'' 在河岸与银行句子中可以分离。


\lecturefigure{slide-014.jpg}{Word embeddings：语义相似性、向量算术与静态表示的局限}{CS25 V5 official deck, p.~14}

\teachervoice{00:06:23--00:07:32，Steven 从“词不是数字”开始，先解释 dense vector，再指出 static embedding 无法区分多义词。这个顺序提醒我们：attention 解决的是上下文交互，不是 tokenization 本身。}

\begin{knowledgebox}{首次使用：tokenization 与 embedding}
Tokenization 把原始字符串切成词、子词或字节 token；embedding 再把 token id 映射为可训练向量。前者决定序列长度与词表边界，后者决定模型初始表示空间。二者经常一起出现，但不是同一个操作。
\end{knowledgebox}

\subsection{Self-attention：Q、K、V 是检索接口}

给定序列矩阵 $X\in\mathbb{R}^{n\times d}$，模型学习
\[
Q=XW_Q,\qquad K=XW_K,\qquad V=XW_V,
\]
并计算 scaled dot-product attention：
\[
\operatorname{Attn}(Q,K,V)=\operatorname{softmax}\!\left(\frac{QK^{\top}}{\sqrt{d_k}}\right)V.
\]
$Q$ 是当前位置提出的问题，$K$ 是每个位置用于匹配的索引，$V$ 是真正被聚合的内容，$d_k$ 是 key 维度；除以 $\sqrt{d_k}$ 用于控制点积尺度。


\lecturefigure{slide-015.jpg}{Transformer 两分钟：self-attention 让 token 软匹配上下文}{CS25 V5 official deck, p.~15}


\lecturefigure{slide-016.jpg}{QKV 图书馆类比：查询、摘要索引与书中内容}{CS25 V5 official deck, p.~16}

\teachervoice{00:07:37--00:08:19，讲者用“找书”解释 Q/K/V：主题需求是 query，书的摘要是 key，书中内容是 value；attention 不是只拿一本书，而是对多个匹配结果做软加权。}

\begin{warningbox}{Attention weight 不是完整解释}
权重说明当前层怎样混合 value，不自动等于因果重要性、语义理由或最终答案证据。残差、MLP、后续层和 value 内容都会改变输出；读 attention map 时应把它当局部计算证据，而不是完整可解释性结论。
\end{warningbox}


\lecturefigure{slide-017.jpg}{Attention map：不同层和头学习不同 token 关系}{CS25 V5 official deck, p.~17}

\subsection{顺序、多头与完整 block}

纯 attention 对 token 排列本身没有位置信号，因此需要 positional encoding。经典正弦形式为
\[
\mathrm{PE}_{(p,2j)}=\sin\!\left(p/10000^{2j/d}\right),\qquad
\mathrm{PE}_{(p,2j+1)}=\cos\!\left(p/10000^{2j/d}\right),
\]
其中 $p$ 是位置，$j$ 是维度索引。现代系统也常用 RoPE 或 learned position；共同目标是让相同 token 在不同位置产生不同交互几何。


\lecturefigure{slide-018.jpg}{Positional encoding：为 permutation-equivariant attention 注入顺序}{CS25 V5 official deck, p.~18}

Multi-head attention 把隐藏维度拆成多个子空间：
\[
\operatorname{MHA}(X)=\operatorname{Concat}(h_1,\ldots,h_H)W_O,
\quad h_r=\operatorname{Attn}(XW_Q^{(r)},XW_K^{(r)},XW_V^{(r)}).
\]
不同 head 可以专注局部语法、长程引用、格式或任务特征；但 head 数增加不等于能力线性增加，真正效果还取决于宽度、数据和训练稳定性。


\lecturefigure{slide-019.jpg}{Multi-head attention：并行关系子空间与层级堆叠}{CS25 V5 official deck, p.~19}


\lecturefigure{slide-020.jpg}{完整 Transformer：attention、MLP、残差与多层组合}{CS25 V5 official deck, p.~20}

\begin{knowledgebox}{深读：一层 Transformer 到底做了哪些资源交换}
Attention 让每个 token 读取其他位置的信息，MLP 则在每个位置独立变换特征；残差把新信息叠加到旧表示，normalization 稳定尺度。训练时必须保存或重算中间 activation，推理时则为历史 token 保存 KV cache。于是同一张 block 图对应三类瓶颈：长序列使 attention 计算和 KV cache 增长，宽 MLP 增加参数与矩阵乘成本，深度增加串行层数和 activation。读架构图不能只数模块，还要把 batch、sequence length、hidden size 和 precision 代入资源账本。所谓“最终 Transformer”只是功能拓扑；FlashAttention、GQA、quantization、parallelism 与 serving scheduler 才决定它在真实硬件上的 latency 和 throughput。图也没有表达训练数据与 optimizer，因此不能从 block 结构直接预测模型最终知识或对齐行为。
\end{knowledgebox}

\teachervoice{00:08:34--00:09:16，Steven 先用 positional encoding 恢复顺序，再用多层和多头解释容量扩张。讲义提醒：``more parameters'' 是表示能力条件，不是数据质量和泛化的替代品。}

\subsection{从 block 到 LLM system}

Transformer 今天横跨语言、视觉、语音、生物和视频，是因为 attention 提供了通用的 token interaction 模板。不同领域真正困难的是 token interface、目标函数与评估约束：图像用 patch，脑成像用 ROI-by-time 序列，agent 还需要 action 和 environment feedback。


\lecturefigure{slide-021.jpg}{Transformers Today：语言、视觉、语音、生物与视频的共同计算模式}{CS25 V5 official deck, p.~21}


\lecturefigure{slide-022.jpg}{Large Language Models：架构、数据、规模、post-training 与应用组成系统}{CS25 V5 official deck, p.~22}

\begin{knowledgebox}{术语消化：从 Transformer 到 LLM system}
\textbf{Base model} 是预训练得到的参数；\textbf{post-training} 用 instruction/preference data 调整行为；\textbf{inference stack} 管理 KV cache、batching 和工具调用；\textbf{product layer} 还包含 retrieval、memory、permissions 与 UI。Slide 22 把这些层压在一张图里。读图时先找每一层的输入输出，再问失败属于哪一层：事实缺失可能是 data 问题，延迟可能是 serving 问题，越权则是 agent policy 问题。图不能证明更大 base model 会自动修复所有上层缺陷。
\end{knowledgebox}

\teachervoice{00:09:20--00:10:10，课堂把 LLM 描述为 scaled-up Transformer 加大规模数据，再经过 post-training 与应用封装。老师强调的不是“架构不重要”，而是架构只是产品能力链条的一段。}

\begin{knowledgebox}{从 next-token loss 到语言模型}
自回归预训练最常见目标是
\[
\mathcal{L}_{\text{NLL}}=-\sum_{t=1}^{T}\log p_{\theta}(x_t\mid x_{<t}).
\]
$x_t$ 是当前 token，$x_{<t}$ 是历史上下文，$\theta$ 是参数。该目标提供统一监督，却不直接规定事实可靠性、工具使用、偏好一致性或长期记忆；这些能力需要数据、post-training 或系统组件共同塑造。
\end{knowledgebox}

\subsection{本章小结}

Transformer 的核心是把离散 token 映射为向量，再通过带位置信号的多头 attention 与 MLP 反复交互。现代 LLM 的能力不能由一个 block 单独解释：训练分布、优化目标、推理时计算、反馈接口和部署环境同样决定结果。下一章转向最常被低估的变量——数据。

